\section{Creativity, Personality, Sex}

\textbf{Julius Evola} has mentioned the philosopher \textbf{Nicholas Berdyaev} several times, including his review of \textit{The Meaning of History}\footnote{\url{https://gornahoor.net/?p=6200}}. But Evola has also mentioned him in the context of the possible recovery of Christianity to tradition. However, Evola ultimately rejected that approach, and I'm not sure if that was meant in the relative sense that Berdyaev's work had little impact in the wider sphere, or in the absolute sense that it was ultimately insufficient in itself. Nevertheless, there are certainly points of contact, such as the emphasis on action and creativity, pre-existence of souls, the priority of the person, not to mention his interest in \textbf{J J Bachofen} and appreciation for \textbf{Friedrich Nietzsche}. Some ideas gleaned from his magnum opus, the Destiny of Man, will give an insight into his thought.

\paragraph{Man as Creative}


Berdyaev sees the essence of human nature as a whole, at least virtually, since there are certainly antinomies and inner conflicts that need to be resolved. To get there, he rejects the evolutionistic explanation, since it makes man a natural creature bound by necessity, rather than free. The Roman view, in his opinion, is limited by seeing man as natural rather than spiritual, who needs grace. The reformed view is worse, seeing man as totally depraved. However, Berdyaev sees in \textbf{Soren Kierkegaard} a ``thinker of genius'' for his psychological insights into anxiety and the paradoxical nature of man. Evola, however, wrote this about that tragic sense of life:

\begin{quotex}
To exalt the romantic, tragic, anxious soul, always in search of new ``truths'', is essentially something of a civilization sick and damaged in its race. Calmness, style, clarity, command, discipline, power, and the Olympic spirit are instead the points of reference for every formation of character and life.
\end{quotex}


Berdyaev then indicates his understanding of the Christian conception of man, which is based on these two ideas:

\begin{enumerate}


\item Man is the image and likeness of God the Creator

\item God became man, the Son of God manifested Himself to us as the God-Man.
\end{enumerate}


Berdyaev next draws out this conclusion from those premises: As the image and likeness of the Creator, man is a creator too and is called to creative cooperation, i.e., he is a creative being. This goes beyond the conceptions mentioned above. This, he emphasizes, is deeper than the idea of man as the ``toolmaker''; in his words:

\begin{quotex}
Man can only be a creative being if he has freedom. There are two elements in human nature, and it is their combination and interaction that constitute man. There is in him the element of primeval, utterly undetermined potential freedom springing from the abyss of non-being, and the element determined by the fact that man is the image and likeness of God, a Divine idea which his freedom may realize or destroy.
\end{quotex}


Assuming this view of man, Berdyaev makes an intriguing point about the Fall, which, he explains, arises from the ``third principle'', i.e.,

\begin{quotex}
uncreated freedom, non-being which is prior to being, the meonic abyss which is neither Creator nor creature ... This is the ultimate mystery behind reality. Endless consequences follow from it. It accounts both for evil and for the creation of what has never existed before. The ethics of creativeness goes back to this primary truth.
\end{quotex}


Yet, for us, this is an elegant restatement of the Hermetic principle of forces. The creature as Destiny, the Creator as Providence, mediated by the Will.

\paragraph{The Person and Pre-Existence}


Another point of contact between Evola and Berdyaev is their understanding of the Person vis-à-vis the individual. For both, the individual is a natural and biological category, whereas the person is spiritual. Berdyaev claims that the

\begin{quotex}
early moral consciousness of mankind is wholly dominated by the mystical power of kinship. Man had to wage a heroic struggle to free himself from it. ... primitive, archaic human morality is entirely communal and traces of it have not completely disappeared among the civilized races of today.
\end{quotex}


Here, unfortunately, Berdyaev is veering close to an evolutionistic understanding of man. Gornahoor in this regard prefers Evola's idea of spiritual races. Rather than indicating an evolution, Berdyaev's observation is better explained by the co-existence of different spiritual races, so that struggle he mentioned is always a here and now battle, and its victory at some point in the past cannot be assumed. Certainly that communal attitude is alive and well in our time as this recent controversial story shows\footnote{\url{http://www.lifesitenews.com/news/msnbc-host-digs-in-heels-over-controversial-psa-claiming-children-belong-to}}. Personhood is an achievement, a difficult achievement in fact, and not common to all.

Berdyaev claims that Christianity freed man from the power of cosmic forces and the blood tie, making the moral life of the individual independent of the tribe or of any collective unit. In his devotion to empire and rejection of the determinism of biological race, Evola likewise sees the person as free of such ties. However, Evola does not see that as exclusively Christian.

Berdyaev concludes his discussion of the person with the idea of the pre-existence of the soul. Obviously, the soul is not a product of the genetic process. Nor is it created in time at the moment of conception, otherwise it could not be free as the third force. Hence, he concludes that it is created by God in eternity, in the spiritual world. Keep in mind that priority can be understood in five ways. Pre-existence in the sense of temporal priority was condemned, but Berdyaev is talking in terms of an ontological priority.

\paragraph{The Masculine and the Feminine}


Along with Kierkegaard, Berdyaev regards \textbf{J J Bachofen} as the other genius in understanding the nature of man. Berdyaev explains why:

\begin{quotex} 
Bachofen discovered the deep primeval layer of human nature, its original connection with the maternal element, the struggle of the masculine solar principle with the feminine tellurgic one, the metaphysics of sex in man. For Bachofen polarity is man's essential characteristic. The cosmic struggle between the sun and the earth, personalism and collectivism, takes place in man.
\end{quotex}


Evola, too, held Bachofen in very high regard, and his adherence to Bachofen's ideas is what makes Evola's thought different from Guenon's in significant ways. Building on Bachofen, Berdyaev understands that man is a sexual being and sexual polarity is characteristic of human nature. This is not merely a biological category but a fundamental quality of man as a whole being. Without mentioning \textbf{Otto Weininger}, Berdyaev's view is amazingly similar. They both understand that the human being combines masculine and feminine elements in different proportions. Here again, the tragic sense of life arises. Although the masculine and the feminine principles seek union, they are also engaged in a cosmic struggle, waging war against each other like deadly enemies.

Bachofen has shown the presence of this struggle in the world. The Sun is the masculine principle: spirit, paternity. The Earth is the feminine principle: matter, flesh, maternity. The Moon is then the ``masculinely feminine intermediary principle''. Bachofen sees history as the result of how these principles meet, interact, and come into conflict within the cosmos.

Berdyaev points out that Bachofen was a Christian. This is important because Berdyaev draws from Bachofen his idea of the person and its separation from blood ties. If original humanity was matriarchal, patriarchy arose with the awakening of the spirit and personality.

Berdyaev next gets pretty alchemical. Following on Plato and \textbf{Jacob Boehme}, man is both personality and cosmos, logos and earth, masculine and feminine:

\begin{quotex}
God's conception of man is a complete, masculinely feminine being, solar and tellurgic, logoic and cosmic at the same time. Only insofar as he is complete is he chaste, wise, and Sophian in his perfect wholeness. As a sexual halved, divided being he is not chaste, or wise, and is doomed to disharmony, to passionate longing and dissatisfaction.
\end{quotex}


Berdyaev refers to the ``horror and curse of sex'', calling it a problem which has never been solved, despite the efforts of the ascetics. Sex is the source of life and of death. Man is sick, wounded, and disharmonious because he has lost his wholeness and integrity. There is also a distinction between procreation and creation. Creativity in man depends on the pre-existing world, since only God creates the world.

Man as person is created by God in eternity, but as individual, he is born, or procreated, in time. The masculine and feminine principles interact and complete each other. Woman inspires man to create. Through creation and procreation man strives to attain the wholeness of his being. Berdyaev sums up his view this way:

\begin{quotex}
The power of sex may be overcome and sublimated. Instead of being generative it may become creative and be a spiritual power. Creativeness is closely allied to sexuality. Sexlessness makes man sterile. A sexless being can neither procreate not create. The moral task is not to destroy the power of sex but to sublimate it, transforming it into a force that creates values. Erotic love is one of such values. Man generates and creates because he is an incomplete being, divided into two and striving for the completeness and wholeness of the androgyne.
\end{quotex}



\flrightit{Posted on 2013-04-16 by Cologero}

\begin{center}* * *\end{center}

\begin{footnotesize}
\begin{sffamily}
\texttt{David on 2013-04-17 at 19:49 said:}

Good articles, is it interesting to see someone else speaking in the terms of Bachofen with somewhat different conclusions. Two things :

1- «man is both personality and cosmos» Here man is taken as humanity, or as masculine ?

2- This is a general question, but most of the time, there is a neoplatonic and dualistic struggle involved (as in Personhood vs Individualhood; that is, soul vs body). Apart from ancient sources (such as the Advaita Vedanta or the Bible), is there any traditional writers or «first line» writers that wrote about a more monistic view of the Personhood within Tradition? Or is neoplatonism the best way to represent the anthropological point of view of Tradition? Thank you

\hfill

\texttt{Cologero on 2013-04-18 at 00:00 said:}

Ironically, both Berdyaev and Evola understand Bachofen in the same way: the solar, virile path is to transcend the matriarchal telluric collective. They just don't see eye to eye on the best path to accomplish that, since Berdyaev is Christian and Evola rejects Christianity. That is why labels are never the whole story.

As for point 1, Berdyaev lists three related dichotomies (man is both personality and cosmos, logos and earth, masculine and feminine). How would you understand that.

Don't forget that there are six schools of Vedanta. The corresponding ``dualist'' school is the Samkhya; it doesn't rule out the non-dual Advaita. So both perspectives are necessary. I think all the writers, at least the ones we are interested in here, have written about the monistic view. This requires the surpassing of the human state. Guido de Giorgio in particular has tried to express that view.

\hfill

\texttt{David on 2013-04-18 at 12:37 said:}

Thank you for the answer. I will try to read more of De Giorgio. As for point one, my question was badly phrased. If man and woman can both have feminine and masculine, when he says man, does he mean man or human ? This may be clear for you, but in french, we dinstinctivly use man or human and so this is not clear for me. Is he writting only in an andropocentric point of view ? Or does this applies to woman as well ? Thank you.

\hfill

\texttt{Cologero on 2013-04-18 at 12:57 said:}

Since the book was originally published in French, you may want to check your library. DE LA DESTINATION DE L'HOMME, Part I, Chapter III, section 3. How do you understand the title? I assume he means women, too, have both characteristics. But more fundamentally, for Berdyaev, there is no ``man'' as such; there are only men and women.

\hfill

\texttt{Jason-Adam on 2013-04-18 at 13:09 said:}

Is it possible to use Bachofen to come to a traditional understanding of the transgender phenomenon we are witnessing in this postmodern world of today where all boundaries and categories are being dissolved ? From some readings of Evola, it seems to me that he was supportive of ``sex changes'' as he did see trannies as ``women trapped in men's bodies.''

I am wondering if the presence of male and female elements in the soul of a man means that things can get reversed and deviations could result, thus causing the formation of ``tranny'' thoughts.

Strangely, Dr Harry Benjamin, one of the pioneers of ``sex changes'' was part of the New York ``fascist'' scene and was a friend of Francis Yockey, Lawrence Dennis and Keith Thompson... ... 

\hfill

\texttt{David on 2013-04-18 at 14:34 said:}

Thank you, I will look into the original version.

\hfill

\texttt{Cologero on 2013-04-19 at 00:15 said:}

J-A, you don't need Bachofen, since formlessness is the hallmark of the kali yuga; it was predicted and it is occurring. The same too for ``deviations'' as you mention. They don't ``cause'' tranny thoughts, since thought is free. However, as some recent posts have shown, thoughts and images arise from all sources. Those that are dwelled upon can lead to acting them out in strange ways.

I had to dig out Dreamers of the Day to catch your references. It calls Dr. Benjamin a ``highly cultured German Jew.'' Apparently, only the ``uncultured'' are uninterested in such deviations. Isn't he advocating a reverse alchemy, turning gold into lead, or more precisely a man into a woman? That is quite remote from the Great Work. I don't see how Yockey, et al, represent any sort of fascist spirit, although it may have been an interesting ``scene'', and it is being reprised now in San Fran.

It is difficult to be counter cultural, or counter current, today, since those values are now the culture and the current. Perhaps you can launch the venerable tradition of right-wing Catholicism as the new hipster scene in NYC, although without all that ``high culture''. Keep me on the evite list for your parties.

\hfill

\texttt{Cologero on 2013-04-21 at 11:08 said:}

I went back to review Evola's thoughts on Bachofen, particularly Ch 27 of Revolt, ``The Civilization of the Mother''. He points out several consequences of the ``dismissal of the virile element'' in the matriarchal cultures. These included:

``Sometimes the priests who felt possessed by the Goddess would go so far as emasculating themselves in order to resemble her and to become transformed into the female type.''

``The Lydian Hercules was dressed as a women''

``Those who participated in some Mysteries often wore women's clothes''

``Priests dressed in women's clothes would keep watch in the sacred woods by some ancient Germanic trees''

Certain cults were celebrated with ``women dressed as men and men dressed as women''

Perhaps Evola ``evolved'' on this issue (and you provided no reference, J-A), but what I quoted is hardly ``supportive'' of your thesis. A different way to look at it is to regard the work of Dr. Benjamin as the symptom of the devolution of the modern world into a gynocractic culture. Her is a further symptom from the same chapter:

``The masculine element and men in general came to be looked down upon as irrelevant, innerly inconsistent, ephemeral, of little value, and as a source of embarrassment.''

In our time, I think it is quite obvious which men in particular are regarded in that way, so I'll leave it to readers to figure it out. A clue: it has nothing to do with ``glorious'' or ``virile'' pagans.

\hfill

\texttt{Jason-Adam on 2013-04-22 at 12:00 said:}

As strange as it might sound, the tradition of right wing Catholicism does live in NYC at St Christopher's SSPX chapel on Lexington Avenue \& 37th Street every Sunday at 1:30 PM under the tutelage of Fr Kevin Robinson, the only priest I have ever met whom every time I am in his prescence I feel humbled to be before a true man of the cloth... ... from personal experience I can say it quite a tough struggle to live tradition in the heart of the city that is alll of the world's decadence but I thrive in a struggle so that is why I chose this life... ... ... ..

You hit upon a noteworthy matter -- Christianity when it first appeared, as I see it, was a truly ``conservative revolution'' in the sense that the pagan Roman society had become utterly and completely corrupt beyond all hope of repair. Guenon goes into this in Crisis of the Modern World, also Fabre D'Olivet had mentioned this as well. Nowadays with society being as immoral as Caligula and Nero (is there any shock over why those 2 men are icons of the fag movements of today ?) the situation allow for the Church to reemerge as the counterculture of sanity again as it was during its first 300 years... ... ... .

While we're on the subject, what is your opinion of Yockey and his ideas ? Dreamer of the Day puzzled me because from mty reading of Imperium and his other works, Yockey struck me as a right wing Catholic but Coogan's biography claims he was a sexual deviant and theosophic-minded occultist... ..I don't trust Coogan the liberal to write truth though... ... 

\hfill

\texttt{Jason-Adam on 2013-04-22 at 12:08 said:}

This is the passage from METAPHSYICS OF SEX that cause me to be believe Evola supported transgenderism and sex changes :

In natural homosexuality or in the predisposition to it, the most straightforward explanation is provided by what we said earlier about the differing levels of sexual development and about the fact that the process of sexual development in its physical and, even more so, in its psychic aspects can be incomplete. In that way, the original bisexual nature is surpassed to a lesser extent than in a ``normal'' human being, the characteristics of one sex not being predominant over those of the other sex to the same extent. Next we must deal with what M. Hirschfeld called the ``intermediate sexual forms''. In cases of this kind (for instance, when a person who is nominally a man is only 60 percent male) it is impossible that the erotic attraction based on the polarity of the sexes in heterosexuality -- which is much stronger the more the man is male and the woman is female -- can also be born between individuals who, according to the birth registry and as regards only the so-called primary sexual characteristics, belong to the same sex, because in actual fact they are ``intermediate forms''. In the case of pederasts, Ulrich said rightly that it is possible to find ``the soul of a woman born in the body of a man''.

\hfill

\texttt{Cologero on 2013-04-22 at 20:07 said:}

This may serve as an explanation, but it not necessarily indicate ``support''. In any case, such a being is ``incomplete'' and an ``intermediate form''. It doesn't follow that any surgical or hormonal reassignment can result in completion.

\hfill

\texttt{Cologero on 2013-04-23 at 07:27 said:}

You may have touched on a nerve, J-A, since Coogan's book soured me on Yockey; perhaps unfairly as you point out. Yet, some, who believe they are part of ``tradition'', are obsessed with the dark, the macabre, the ugly, the perverse, the deviant, bondage, domination, etc., in short, everything that is the opposite of solarity. And their footprint on the Internet is not insignificant.

\hfill

\texttt{Jason-Adam on 2013-04-23 at 13:39 said:}

I fully agree, the worldview of Gothic rock has no place in the life of an aristocrat. Light shall triumph over the darkness.

\hfill

\texttt{Jason-Adam on 2013-04-23 at 13:42 said:}

Indeed, a ``transwoman'' is simply a man who has mutilated his body, as the Church teaches us -- so the question I have regarding fags, trannies and other deviants is -- should they be ``saved'' via deprogramming as many Christians attempt to do, with poor results usually, or should we just allow them to exist as the dregs of society as Evola -- and also the Hindus and Buddhists -- counsels ?

\end{sffamily}
\end{footnotesize}

